%\vspace{1.25in}

\question
\begin{multicols}{2}
Use the phylogenetic tree at right to answer the following questions.

\begin{parts}
\part[3]
What does the branch labeled ``1'' represent on this tree? What is the proper term for this branch?

\begin{minipage}[t]{0.9\columnwidth}%
\AnswerBox{0.15\textheight}{%
The root (1 pt) represents the common ancestor of all groups shown on the phylogeny (2 pts). }
\end{minipage}

\part[2]
Name the basal taxon for this phylogeny.  

\begin{minipage}[t]{0.9\columnwidth}%
\AnswerBox{0.03\textheight}{
Porifera.}
\end{minipage}

\part[5]
Are Platyhelminthes more closely related to Molluska or to Rotifera. Tell how you know.

\begin{minipage}[t]{0.9\columnwidth}%
\AnswerBox{0.22\textheight}{
Platyhelminthes are more closely related to Molluska (2 pts) because they share a more recent common ancestor with mollusks than rotifers (2 pts).}
\end{minipage}

\part[5]
The taxon Protostomia includes Rotifera, Platyhelminthes, Annelida, Molluska, Nematoda, and Arthropoda. Are the Protostomia a monophyletic group? Tell how you know.

\begin{minipage}[t]{0.9\columnwidth}%
\AnswerBox{0.22\textheight}{%
Protosomes are  monophyletic (2 pts). The groups contain some but not all of the descendants of the common ancestor (2 pts).}
\end{minipage}


%\part[5]
%Identify a pair of sister taxa. Tell why they are sister taxa.

%\begin{solution}
%	They share a most recent common ancestor.
%\end{solution}

\end{parts}
\columnbreak
\begin{center}
\includegraphics[width=\columnwidth]{exam2_protostome_phylogeny}
\end{center}
\end{multicols}